\documentclass[11pt,twoside,a4paper]{article}
% http://www-h.eng.cam.ac.uk/help/tpl/textprocessing/latex_maths+pix/node6.html symboles de math
% http://fr.wikibooks.org/wiki/Programmation_LaTeX Programmation latex (wikibook)
%=========================== En-Tete =================================
%--- Insertion de paquetages (optionnel) ---
\usepackage[french]{babel}   % pour dire que le texte est en fran{\'e}ais
\usepackage{a4}	             % pour la taille   
\usepackage[T1]{fontenc}     % pour les font postscript
\usepackage{epsfig}          % pour gerer les images
%\usepackage{psfig}
\usepackage{amsmath, amsthm} % tres bon mode mathematique
\usepackage{amsfonts,amssymb}% permet la definition des ensembles
\usepackage{float}           % pour le placement des figure
\usepackage{verbatim}

\usepackage{longtable} % pour les tableaux de plusieurs pages

\usepackage[table]{xcolor} % couleur de fond des cellules de tableaux

\usepackage{lastpage}

\usepackage{multirow}

\usepackage{multicol} % pour {\'e}crire dans certaines zones en colonnes : \begin{multicols}{nb colonnes}...\end{multicols}

%% https://texblog.org/2011/02/26/generating-dummy-textblindtext-with-latex-for-testing/
%% https://fr.sharelatex.com/learn/Multiple_columns
%% \usepackage{blindtext}
%% \usepackage{lipsum}
\usepackage{wrapfig}

% \usepackage[top=1.5cm, bottom=1.5cm, left=1.5cm, right=1.5cm]{geometry}
% gauche, haut, droite, bas, entete, ente2txt, pied, txt2pied
\usepackage{vmargin}
\setmarginsrb{1.0cm}{1.0cm}{1.0cm}{1.0cm}{15pt}{3pt}{57pt}{3pt}

\usepackage{lscape} % changement orientation page
\usepackage{pdflscape}
% --- style de page (pour les en-tete) ---
% \pagestyle{headings}
\pagestyle{empty}

% % % en-tete et pieds de page configurables : fancyhdr.sty

% http://www.trustonme.net/didactels/250.html

% http://ww3.ac-poitiers.fr/math/tex/pratique/entete/entete.htm
% http://www.ctan.org/tex-archive/macros/latex/contrib/fancyhdr/fancyhdr.pdf
% \usepackage{fancyhdr}
% \pagestyle{fancy}
% % \newcommand{\chaptermark}[1]{\markboth{#1}{}}
% % \newcommand{\sectionmark}[1]{\markright{\thesection\ #1}}
% \fancyhf{}
% \fancyhead[LE,RO]{\bfseries\thepage}
% \fancyhead[LO]{\bfseries\rightmark}
% \fancyhead[RE]{\bfseries\leftmark}
% \fancyfoot[LE]{\thepage /\pageref{LastPage} \hfill
	% TITLE
% \hfill \includegraphics[width=0.5cm]{img/logo_glider.png} }
% \fancyfoot[RO]{\includegraphics[width=0.5cm]{img/logo_glider.png} \hfill
	% TITLE
% \hfill \thepage /\pageref{LastPage}}
% \renewcommand{\headrulewidth}{0.5pt}
% \renewcommand{\footrulewidth}{0.5pt}
% \addtolength{\headheight}{0.5pt}
% \fancypagestyle{plain}{
	% \fancyhead{}
	% \renewcommand{\headrulewidth}{0pt}
% }

%--- Definitions de nouvelles commandes ---
\newcommand{\N}{\mathbb{N}} % les entiers naturels

%============================= Corps =================================
\begin{document}
\begin{landscape}

% \setcounter{page}{0}
% \thispagestyle{empty}

\begin{multicols}{2}
	\section*{Memento REST}
	
	\begin{itemize}
		\item Services Web
		\begin{itemize}
			\item Protocoles standardis{\'e}s / ouverts (en g{\'e}n{\'e}ral HTTP)
			\item Deux familles : SOAP (orient{\'e} services) et REST (orient{\'e} ressources)
			\item REpresentational State Transfer, chapitre 5 th{\`e}se de Roy Fielding (2000), syst{\`e}mes distribu{\'e}s, \texttt{http://opikanoba.org/tr/fielding/rest}
			\item Architecture, ensemble de contraintes... 
			\item Client / serveur ; Stateless ; interface uniforme (identification, repr{\'e}sentation, auto-descriptif, HATEOAS / Hypermedia) ; couches ; code {\`a} la demande (optionnel)
		\end{itemize}
		\item Client / serveur : s{\'e}pare IHM et stockage ({\'e}volution s{\'e}par{\'e}e, meilleure portabilit{\'e} et possibilit{\'e}s mont{\'e}e en charge...)
		\item Stateless : requ{\^e}te auto-portante ; pas de contexte sur le serveur (c{\^o}t{\'e} client) ; am{\'e}liore visibilit{\'e}
		\item Donn{\'e}es r{\'e}ponses marqu{\'e}es pour la mise en cache
		\item Interfaces uniformes (standardis{\'e}) : cf. SQL / SGBDR (schema.table.colonne) : SELECT, INSERT, UPDATE et DELETE
		\item Contraintes d'interface : identification uniforme (URI : Uniform Resource Identifier)...
		\item Repr{\'e}sentations (seule la repr{\'e}sentation de l'{\'e}tat transite)
		\item Hyperm{\'e}dia comme moteur de l'{\'e}tat de l'application (HATEOAS : Hypertext As The Engine Of Application State) : facilit{\'e}, guidage utilisateur, agent utilisateur stocke l'information, pas de contraintes pr{\'e}-d{\'e}finies...
		\item ...
		\item HTTP / URI = Uniform Resource Identifier (RFC 3986) \texttt{https://tools.ietf.org/html/rfc3986}
		\begin{itemize}
			\item URL : Uniform Resource Locator (identifie une ressource sur un r{\'e}seau)
			\item URN : Uniform Resource Name (identifie un nom de ressource)
			\item Schema + Chemin d'acc{\`e}s
		\end{itemize}
		
		\item Requ{\^e}te : Structure textuelle normalis{\'e}e
		\begin{itemize}
			\item Ligne de requ{\^e}te : <m{\'e}thode> <URI> <version-HTTP> CRLF
			\item Lignes d'en-t{\^e}te : *(<champs>: <valeur> CRLF)
			\item Ligne vide : CRLF
			\item Corps (optionnel) : [*OCTET]
		\end{itemize}
	\end{itemize}
	
	~\\
	
	\columnbreak
	
	\begin{itemize}

		\item R{\'e}ponse : Structure assez similaire {\`a} la requ{\^e}te
		\begin{itemize}
			\item Ligne de statut : <version-HTTP> <statut> <message> CRLF
			\item Lignes d'en-t{\^e}te : *(<champs>: <valeur> CRLF)
			\item Ligne vide : CRLF
			\item Corps (optionnel) : [*OCTET]
		\end{itemize}
		\item De nombreux champs d'en-t{\^e}te sont pr{\'e}d{\'e}finis
		\begin{itemize}
			\item General header : les champs communs (Cache-Control, Date, ...)
			\item Request header : les champs sp{\'e}cifiques de requ{\^e}te (Host, Accept, ...)
			\item Response header : les champs sp{\'e}cifiques de r{\'e}ponse (Age, ETag, ...)
			\item Entity header : les champs sp{\'e}cifiques des entit{\'e}s (Content-Encoding, ...)
		\end{itemize}
		\item Sp{\'e}cifications : http://www.w3.org/Protocols/
		\begin{itemize}
			\item HTTP 1.1 : https://tools.ietf.org/html/rfc2616
			\item HTTP 1.1bis : reformulation dans plusieurs documents (RFC7230, …, RFC7237)
			\item HTTP 2 : RFC7540 publi{\'e}e en mai 2015
		\end{itemize}
		\item 8 m{\'e}thodes de base
		\begin{itemize}
			\item GET : obtenir une repr{\'e}sentation d'une ressource
			\item HEAD : demander des informations sur une ressource
			\item POST : transmettre des donn{\'e}es {\`a} une ressource pour traitement
			\item PUT : remplacer ou cr{\'e}er une ressource (dont on conna{\^i}t l'URI)
			\item DELETE : supprimer une ressource
			\item CONNECT : pour utiliser un Proxy comme un tunnel de communication
			\item OPTIONS : demander les options de communications (ressource ou serveur)
			\item TRACE : demander au serveur de retourner ce qu'il a re\c{c}u (echo pour diagnostic)
		\end{itemize}
		\item Codes de r{\'e}ponse HTTP : 
	\end{itemize}
	
	\begin{wraptable}{ht}{0.8\linewidth}
		\centering
		\begin{tabular}{|l|l|l|}
		\hline
			Code	&					&	Signification										\\
			\hline
			1xx		&	Information		&	messages d'information (non utilis{\'e})			\\
			2xx		&	Succ{\`e}s		&	la requ{\^e}te s'est d{\'e}roul{\'e}e correctement	\\
			3xx		&	Redirection		&	indique une redirection								\\
			4xx		&	Erreur client	&	{\`a} l'origine du client							\\
			5xx		&	Erreur serveur	&	{\`a} l'origine du serveur							\\
			\hline
		\end{tabular}
		\caption{Table r{\'e}sumant les codes HTTP}
		\label{table:httpcodes}
	\end{wraptable}
	
	%% \vfill
	
	...
\end{multicols}

\clearpage

\end{landscape}
\end{document}
